\section*{Description du problème}

Nous proposons de modéliser et résoudre un problème réel de planification de maintenance ferroviaire sur 24h,
inspiré directement des opérations au centre de maintenance SNCF de Paris.
Ce problème nous a été rapporté par des connaissances ayant travaillé sur site,
où la planification quotidienne est actuellement réalisée manuellement via Excel.
Cette approche manuelle est extrêmement chronophage :
les planificateurs passent plusieurs heures chaque jour à construire et ajuster le planning,
souvent en début de journée et à chaque modification imprévue.

L'aspect innovant de ce projet réside dans le fait que la SNCF n'a jamais formalisé ce problème comme un problème d'optimisation combinatoire.
Notre démarche de formalisation et de résolution par programmation par contraintes constitue
donc une contribution potentiellement très utile pour l'industrie ferroviaire, en transformant un processus manuel et
chronophage en un système automatisé capable de générer rapidement des plannings optimisés.

Le planificateur doit quotidiennement organiser les interventions de maintenance sur les trains disponibles en respectant de multiples contraintes :

\begin{itemize}
	\item \textbf{Fenêtres temporelles :}
	      Chaque train a une fenêtre de disponibilité (temps entre son arrivée à l'atelier et son prochain départ).
	      Ces fenêtres varient fortement (de 2h à 24h) et peuvent changer.

	\item \textbf{Équipes spécialisées :}
	      Les opérations nécessitent différentes équipes (mécanique, électricité, confort) avec des compétences spécifiques.
	      Certaines équipes ne peuvent pas se remplacer et les effectifs sont limités.

	\item \textbf{Capacité d'atelier :}
	      Nombre limité de voies de garage, certains travaux nécessitent des conditions spéciales (alimentation électrique, équipements particuliers).

	\item \textbf{Priorités :}
	      Certaines opérations ont des dates butoirs obligatoires (maintenance préventive régulière), d'autres sont urgentes suite à des pannes.

	\item \textbf{Modifications :}
	      Des travaux non terminés doivent être reportés, des trains peuvent être ajoutés ou retirés, nécessitant une replanification rapide.
\end{itemize}

Ce problème combine ordonnancement de tâches avec affectation de ressources multiples et fenêtres temporelles contraintes.
